% !PRINTED 146
\weberpage{146}
%<<CONT
\weq{8}{\begin{gathered}
 C_0\alpha_1^{n-1}+C_1\alpha_1^{n-2}+\cdots+C_{n-2}\alpha_1+(C_{n-1}-S)=0\\
 C_0\alpha_2^{n-1}+C_1\alpha_2^{n-2}+\cdots+C_{n-2}\alpha_2+(C_{n-1}-S)=0\\
 \makebox[30em]{\dotfill}\\
 C_0\alpha_n^{n-1}+C_1\alpha_n^{n-2}+\cdots+C_{n-2}\alpha_n+(C_{n-1}-S)=0.
\end{gathered}}
If we regard this system as a system of homogeneous linear equations with the $n$ unknowns
\[
 C_0,C_1\ldots C_{n-2},C_{n-1}-S,
\]
then its determinant,
\[
\begin{vmatrix}
 \alpha_1^{n-1},&\alpha_1^{n-2}&\ldots&\alpha_1,&1\\
 \alpha_2^{n-1},&\alpha_2^{n-2}&\ldots&\alpha_2,&1\\
 \multicolumn{5}{c}{\makebox[12em]{\dotfill}}\\
 \alpha_n^{n-1},&\alpha_n^{n-2}&\ldots&\alpha_n,&1
\end{vmatrix},
\]
which by \S\,22, formula (12), equals the product of all differences $\alpha_1-\alpha_2,\alpha_1-\alpha_3\ldots$, is different from zero, and consequently, by Theorem II in \S\,24,
\[
 C_0=0,\quad C_1=0\ldots\quad C_{n-2}=0
\]
and
\weq{9}{S=C_{n-1},}
which contains the fundamental theorem to be proved.

The same result is also obtained by Theorem II, \S\,30; for since the function of degree $(n-1)$ in $x$
\[
 C_0x^{n-1}+C_1x^{n-2}+\cdots+C_{n-2}x+C_{n-1}-S
\]
vanishes for $x=\alpha_1,\alpha_2\ldots\alpha_n$, that is, for more than $n-1$ values, all its coefficients must be zero.

The proof which we have given here for the fundamental theorem, following \gesperrt{Cauchy},\footnote[1]{\gesperrt{Cauchy}, Exercices de mathématiques, $4^{\text{ème}}$ année.} also provides a means of actually computing, in particular cases, the expression of a symmetric function in terms of the fundamental functions. The other known proofs of the theorem offer the same possibility. We shall give one more proof here, which often leads to a simpler method of computation\footnote[2]{\gesperrt{Waring}, Meditationes algebraicae. \gesperrt{Gauss}, Demonstratio nova altera theorematis omnem functionem algebraicam rationalem integram unius variabilis in factores reales primi vel secundi gradus resolvi posse. Werke, vol.\ III, p.\ 36.}.
